\section{Quick Start}\label{sec:quick-start}

\begin{enumerate}
  \item Connect an oscillator or other audio source to the red input. Keep
  your existing connection to the mixer or audio output if you want to hear it.
  \item Leave the input gain at $0\,\mathrm{dB}$ and ensure \textbf{Run} is lit.
  A new module uses a Flattop window, Length 2048, and Hop
  $30\,\mathrm{ms}$, with no time or frequency smoothing.
  \item For an amplitude comparison, set \textbf{Slope} to
  $0\,\mathrm{dB/oct}$. The default $+4.5\,\mathrm{dB/oct}$ visually boosts
  frequencies above $1\,\mathrm{kHz}$ and reduces those below it.
  \item Connect another signal to the green, blue, or yellow input to
  compare spectra. Use equal input gains for a level comparison.
  \item Hover over the plot to show a crosshair, frequency, musical note
  and tuning offset, and amplitude at the cursor position. The readout
  describes the cursor; it does not automatically find a spectral peak.
  \item Press \textbf{Run} to stop capturing new samples. Press it again
  to resume. The retained audio can still be reanalyzed while capture is stopped.
\end{enumerate}

\paragraph{Reading levels}
Fourier displays spectral magnitudes, not total signal RMS, loudness, or
sound pressure level. Its nominal amplitude reference is a
$5\,\mathrm{V}$ peak sine wave (approximately $0\,\mathrm{dB}$ or
$100\,\%$ with unity gain, zero slope, and no smoothing). Window choice,
frequency alignment with FFT bins, and smoothing affect the height and width
of a peak. The display ceiling is about $+12\,\mathrm{dB}$; a trace reaching
that ceiling does not mean Fourier has clipped your audio path.

\paragraph{Choosing settings}
Use a longer \textbf{Length} to separate nearby frequencies, especially in
the bass. Use a shorter \textbf{Hop} and little or no \textbf{Average} to
follow changes more closely. Increase \textbf{Average} for a steadier trace
and \textbf{Smooth} for a broader view of tonal balance. The factory presets
provide examples of these choices.
